\documentclass[10pt]{article}
\usepackage[a4paper,margin=0.55in]{geometry}
\usepackage{amsmath,amssymb,amsfonts,mathtools}
\usepackage[dvipsnames]{xcolor}
\usepackage{array,booktabs,float}
\usepackage[colorlinks=true,linkcolor=blue]{hyperref}
\allowdisplaybreaks
\numberwithin{equation}{section}
\setlength{\jot}{4pt}
\setlength{\emergencystretch}{3em}

\title{Group I, Row V, part 2b: UV term (dim.\ reg., $\overline{\rm MS}$) $+$ finite term,\\ with $d_\perp$ kept general}
\date{}

\begin{document}
\maketitle
\noindent This note does the following:
\begin{itemize}
\item It writes your part 2b as a UV-divergent term plus a finite term.
\item The divergence at $x\to z$ is regulated in $d=2-\epsilon$ transverse dimensions in $\overline{\rm MS}$, as in Group I Row II and Row IV.
\item The polarization trace $\delta^{ii}=d_\perp$ of your first bracket is kept general. Setting $d_\perp=2-\epsilon$ reproduces the result in \texttt{Part2b\_MSbar.pdf} and \texttt{Group\_I\_corrected.pdf}.
\item Nothing is expanded in $\Lambda/k^+$.
\end{itemize}

%==========================================================================================
\section{Result}\label{sec:result}
\begin{equation}
\boxed{\frac{d^{3}N_{2b}}{d^{3}k}=\frac{d^{3}N_{2b}}{d^{3}k}\bigg|_{\rm UV}+\frac{d^{3}N_{2b}}{d^{3}k}\bigg|_{\rm finite}\qquad(+\ {\rm c.c.})}
\end{equation}

\paragraph{First term: UV divergent.}
\begin{align}
\frac{d^{3}N_{2b}}{d^{3}k}\bigg|_{\rm UV}={}&\frac{1}{(2\pi)^{3}}\,\frac{g^{4}N_{c}}{16\pi^{4}}\,\frac{1}{k^{+}}\int_{y,z,x',w'}e^{-ik\cdot(w'-z)}\,\frac{(x'-w')\cdot(y-z)}{(x'-w')^{2}\,(y-z)^{2}}\notag\\
&\times\Big[U^{ab'}(x')-U^{ab'}(w')\Big]\Big[U^{af}(y)-U^{af}(z)\Big]\rho^{b'}(x')\,\rho^{f}(y)\notag\\
&\times\bigg\{\Big[2\log\frac{\vee}{\Lambda}-2\log\frac{\vee}{k^{+}}+2\log\Big(1-\frac{\Lambda}{k^{+}}\Big)-2+\frac{4\Lambda}{k^{+}}\Big]\Big[\frac{2}{\epsilon}+2\gamma_{E}+\log\frac{(y-z)^{2}\mu^{2}_{\overline{MS}}}{4}\Big]\notag\\
&\qquad+\frac{d_\perp}{2}\Big[\frac16-\frac{\Lambda^{2}}{k^{+2}}+\frac{2\Lambda^{3}}{3k^{+3}}\Big]\Big[\frac{2}{\epsilon}+2\gamma_{E}+\log\frac{(y-z)^{2}\mu^{2}_{\overline{MS}}}{4}+1\Big]\bigg\}.
\label{eq:UV}
\end{align}
\begin{itemize}
\item The pole is UV: $\epsilon>0$.
\item The ``$+1$'' in the second line comes from $\frac{d_\perp}{d}=\frac{d_\perp}{2}\big(1+\frac\epsilon2\big)$ multiplying $\frac2\epsilon$. Here $d=2-\epsilon$ is the dimension of the $x$ integral.
\item The $N_c$ of your colour factor $[N_cU(y)-N_cU(x)]$ is in the prefactor.
\end{itemize}

\paragraph{With $d_\perp=2-\epsilon$} (conventional dimensional regularization, as in Row II and Row IV), $\frac{d_\perp}{2}\big[\frac2\epsilon+\dots+1\big]=\frac2\epsilon+\dots$, and $-2+\frac16=-\frac{11}{6}$:
\begin{align}
\frac{d^{3}N_{2b}}{d^{3}k}\bigg|_{\rm UV}={}&\frac{1}{(2\pi)^{3}}\,\frac{g^{4}N_{c}}{16\pi^{4}}\,\frac{1}{k^{+}}\Big[2\log\frac{\vee}{\Lambda}-2\log\frac{\vee}{k^{+}}+2\log\Big(1-\frac{\Lambda}{k^{+}}\Big)-\frac{11}{6}+\frac{4\Lambda}{k^{+}}-\frac{\Lambda^{2}}{k^{+2}}+\frac{2\Lambda^{3}}{3k^{+3}}\Big]\notag\\
&\times\int_{y,z,x',w'}e^{-ik\cdot(w'-z)}\Big[\frac{2}{\epsilon}+2\gamma_{E}+\log\frac{(y-z)^{2}\mu^{2}_{\overline{MS}}}{4}\Big]\frac{(x'-w')\cdot(y-z)}{(x'-w')^{2}\,(y-z)^{2}}\notag\\
&\times\Big[U^{ab'}(x')-U^{ab'}(w')\Big]\Big[U^{af}(y)-U^{af}(z)\Big]\rho^{b'}(x')\,\rho^{f}(y)
\label{eq:UVcdr}
\end{align}
(\texttt{sympy}: difference from Eq.~\eqref{eq:UV} at $d_\perp=2-\epsilon$ $=$ @@CDR@@).

\paragraph{$\overline{\rm MS}$.} The factor $\big(\frac{\mu^2e^{\gamma_E}}{4\pi}\big)^{\epsilon/2}$ of the $x$ integral is already in $\mu_{\overline{MS}}$, so $\overline{\rm MS}$ removes only $\frac2\epsilon$:
\begin{equation}
\frac{2}{\epsilon}+2\gamma_{E}+\log\frac{(y-z)^{2}\mu^{2}_{\overline{MS}}}{4}\ \longrightarrow\ 2\gamma_E+\log\frac{(y-z)^{2}\mu^{2}_{\overline{MS}}}{4}=\log\frac{(y-z)^{2}\mu^{2}_{\overline{MS}}e^{2\gamma_E}}{4}.
\end{equation}

\paragraph{Second term: finite} ($d=2$; $d_\perp$ as in your first bracket):
\begin{align}
\frac{d^{3}N_{2b}}{d^{3}k}\bigg|_{\rm finite}={}&\frac{1}{(2\pi)^{3}}\,\frac{g^{4}}{8\pi^{4}}\,\frac{1}{k^{+}}\int_{\Lambda}^{k^{+}-\Lambda}\frac{dp^{+}}{2\pi}\int_{x,y,z,x',w'}e^{-ik\cdot(w'-z)}\,\Bigg\{\notag\\
&e^{-ik\cdot(z-x)\frac{p^{+}}{k^{+}}}\,\frac{1}{\left[p^{+}(y-x)^{2}+(k^{+}-p^{+})(y-z)^{2}\right]}\,\frac{(x'-w')^{i}\,(z-x)^{m}}{(x'-w')^{2}\,(z-x)^{2}}\notag\\
&\times\Bigg[\frac{\delta^{im}\,d_{\perp}}{2\,(x-z)^{2}}\left[(y-x)^{2}-(y-z)^{2}\right]\frac{p^{+}(k^{+}-p^{+})}{k^{+2}}
-\frac{p^{+}}{k^{+}-p^{+}}\left[\frac{(y-x)^{i}(x-z)^{m}}{(x-z)^{2}}-\frac{(y-x)^{i}(y-z)^{m}}{2(y-z)^{2}}\right]\notag\\
&\qquad+\frac{(k^{+}-p^{+})\,\delta_{im}}{k^{+}}\left[\frac{(y-z)\cdot(x-z)}{(x-z)^{2}}+\frac{(y-x)\cdot(y-z)}{2(y-x)^{2}}-\frac{(y-x)^{2}-(y-z)^{2}}{2(x-z)^{2}}\right]\notag\\
&\qquad-\frac{k^{+}-p^{+}}{p^{+}}\left[\frac{(y-z)^{i}(x-z)^{m}}{(x-z)^{2}}+\frac{(y-x)^{m}(y-z)^{i}}{2(y-x)^{2}}\right]\notag\\
&\qquad+\frac{p^{+}\,\delta_{im}}{k^{+}}\left[\frac{(y-x)\cdot(x-z)}{(x-z)^{2}}-\frac{(y-x)\cdot(y-z)}{2(y-z)^{2}}-\frac{(y-x)^{2}-(y-z)^{2}}{2(x-z)^{2}}\right]\notag\\
&\qquad-\left[\frac{(y-x)^{m}(x-z)^{i}}{(x-z)^{2}}-\frac{(y-x)^{m}(y-z)^{i}}{2(y-z)^{2}}+\frac{(y-z)^{m}(x-z)^{i}}{(x-z)^{2}}+\frac{(y-x)^{i}(y-z)^{m}}{2(y-x)^{2}}\right]\Bigg]\notag\\
&\times\Big[U^{ab'}(x')-U^{ab'}(w')\Big]\Big[N_{c}\,U^{af}(y)-N_{c}\,U^{af}(x)\Big]\rho^{b'}(x')\,\rho^{f}(y)\notag\\
&-\frac{1}{k^{+}}\left[\frac{d_\perp}{2}\frac{p^{+}(k^{+}-p^{+})}{k^{+2}}+\frac{p^{+}}{k^{+}-p^{+}}+\frac{k^{+}-p^{+}}{p^{+}}\right]
\frac{(x'-w')\cdot(y-z)}{(x'-w')^{2}\,(y-z)^{2}}\,\frac{(y-x)\cdot(y-z)}{(y-x)^{2}\,(x-z)^{2}}\notag\\
&\times\Big[U^{ab'}(x')-U^{ab'}(w')\Big]\Big[N_{c}\,U^{af}(y)-N_{c}\,U^{af}(z)\Big]\rho^{b'}(x')\,\rho^{f}(y)\Bigg\}.
\label{eq:fin}
\end{align}
The first six lines inside the braces are your 2b, unchanged. The last two lines subtract the counterterm: no phase $e^{-ik\cdot(z-x)p^+/k^+}$, colour $[N_cU(y)-N_cU(z)]$, kernel $\frac{(y-x)\cdot(y-z)}{(y-x)^2(x-z)^2}$. The finite term is UV finite and is evaluated in two dimensions. Its $p^+$ integral is over $\Lambda<p^+<k^+-\Lambda$ with no plus prescription. For $d_\perp=2$ it is Eq.~(1.5) of \texttt{Part2b\_MSbar.pdf}.

%==========================================================================================
\section{Derivation}\label{sec:der}
Notation: $\xi\equiv\frac{p^+}{k^+}$, $\mathbf r\equiv x-z$, $Y\equiv y-z$, and $K(a,b,c)\equiv\frac{(a-c)\cdot(b-c)}{(a-c)^2(b-c)^2}$ (the third argument is integrated over).

\subsection*{Step 1: where it diverges}
At $x\to z$ the integrand behaves like $\frac{1}{(x-z)^2}$. The colour factor $N_c[U(y)-U(x)]\to N_c[U(y)-U(z)]$ does not vanish there, so $\int d^2x$ diverges logarithmically.

The other regions are finite:
\begin{itemize}
\item $y\to x$ and $y\to z$ are at most $\frac1{|y-x|}$ or $\frac1{|y-z|}$;
\item at $x,y,z$ all together the integrand scales as $\rho^{-3}$, against $d^2x\,d^2y\sim\rho^3d\rho$;
\item at large $|x-z|$ the $x$ integral converges through the phase $e^{-ik\cdot(z-x)p^+/k^+}$.
\end{itemize}

\subsection*{Step 2: the short-distance limit (with \texorpdfstring{$d_\perp\neq d$}{d_perp not equal to d})}
Set $z=0$ and $x=\mathbf r\to0$. Then:
\begin{itemize}
\item the phase $\to1$;
\item $\frac{1}{p^+(y-x)^2+(k^+-p^+)(y-z)^2}\to\frac{1}{k^+(y-z)^2}$;
\item $\frac{(z-x)^m}{(z-x)^2}=-\frac{r^m}{r^2}$.
\end{itemize}
Average over the directions of $\mathbf r$ in $d$ dimensions, using $\langle r^ir^j\rangle=\frac{r^2}{d}\delta^{ij}$. Call your six brackets $B_1,\dots,B_6$; they give:
\begin{center}\small
\begin{tabular}{@{}cccc@{}}\toprule
term & $p^+$ factor & leading term of $-\frac{r^m}{r^2}B^{im}$ & angular average\\\midrule
$B_1$ & $\xi(1-\xi)$ & $d_\perp\frac{r^i(Y\cdot r)}{r^4}-\frac{d_\perp}{2}\frac{r^i}{r^2}$ & $\frac{d_\perp}{d}\frac{Y^i}{r^2}$\\
$B_2$ & $-\frac{\xi}{1-\xi}$ & $-\frac{Y^i}{r^2}+\frac{r^i}{r^2}+\frac{(Y\cdot r)Y^i}{2Y^2r^2}$ & $-\frac{Y^i}{r^2}$\\
$B_3$ & $1-\xi$ & $-2\frac{r^i(Y\cdot r)}{r^4}$ & $-\frac2d\frac{Y^i}{r^2}$\\
$B_4$ & $-\frac{1-\xi}{\xi}$ & $-\frac{Y^i}{r^2}-\frac{(Y\cdot r)Y^i}{2Y^2r^2}$ & $-\frac{Y^i}{r^2}$\\
$B_5$ & $\xi$ & $-2\frac{r^i(Y\cdot r)}{r^4}+2\frac{r^i}{r^2}$ & $-\frac2d\frac{Y^i}{r^2}$\\
$B_6$ & $-1$ & $-2\frac{r^i(Y\cdot r)}{r^4}+\frac{r^i}{r^2}$ & $-\frac2d\frac{Y^i}{r^2}$\\\bottomrule
\end{tabular}
\end{center}
With the $p^+$ factors:
\begin{equation}
\frac{d_\perp}{d}\,\xi(1-\xi)+\frac{\xi}{1-\xi}-\frac{2}{d}(1-\xi)+\frac{1-\xi}{\xi}-\frac2d\,\xi+\frac2d=\frac{d_\perp}{d}\,\xi(1-\xi)+\frac{\xi}{1-\xi}+\frac{1-\xi}{\xi}.
\label{eq:coef}
\end{equation}
The $\frac1d$ terms of $B_3,B_5,B_6$ cancel. Only $B_1$ keeps $\frac{d_\perp}{d}$. Numerically, for $d=2,3$ and $d_\perp=1,2,3,5$: rel.\ @@SD@@. So, at $x\to z$, your 2b is
\begin{align*}
&\frac{1}{(2\pi)^{3}}\frac{g^{4}}{8\pi^{4}}\frac{1}{k^{+}}\frac{dp^+}{2\pi}\,e^{-ik\cdot(w'-z)}\,\frac{1}{k^{+}}\Big[\frac{d_\perp}{d}\frac{p^{+}(k^{+}-p^{+})}{k^{+2}}+\frac{p^{+}}{k^{+}-p^{+}}+\frac{k^{+}-p^{+}}{p^{+}}\Big]\\
&\times\frac{(x'-w')\cdot(y-z)}{(x'-w')^{2}(y-z)^{2}}\,\frac{1}{(x-z)^{2}}\Big[U^{ab'}(x')-U^{ab'}(w')\Big]\Big[N_cU^{af}(y)-N_cU^{af}(z)\Big]\rho^{b'}(x')\rho^{f}(y)
\end{align*}
plus terms with zero angular average and $O(\frac1{|x-z|})$.

\subsection*{Step 3: the counterterm}
Replace $\frac{1}{(x-z)^2}$ by
\begin{equation}
K(z,z,x)-K(y,z,x)=\frac{(y-x)\cdot(y-z)}{(y-x)^{2}(x-z)^{2}} .
\end{equation}
It has the following properties:
\begin{itemize}
\item the same $\frac{1}{(x-z)^2}$ at $x\to z$, up to a term with zero angular average;
\item it falls like $|x-z|^{-3}$ at large $|x-z|$;
\item in two dimensions its angular average is exactly $\frac{\theta(|y-z|-|x-z|)}{(x-z)^2}$ (checked: @@SAV@@).
\end{itemize}
The counterterm CT is Step 2 with this kernel, integrated over $\Lambda<p^+<k^+-\Lambda$. Then
\begin{itemize}
\item $\frac{d^3N_{2b}}{d^3k}=(2b-{\rm CT})+{\rm CT}$;
\item CT, done in $d$ dimensions, is the UV term;
\item $2b-{\rm CT}$ at $d=2$ is the finite term. In it the $\frac{d_\perp}{d}$ becomes $\frac{d_\perp}{2}$, Eq.~\eqref{eq:fin}.
\end{itemize}

\subsection*{Step 4: the \texorpdfstring{$x$}{x} integral in \texorpdfstring{$d=2-\epsilon$}{d=2-eps}, \texorpdfstring{$\overline{\rm MS}$}{MS-bar}}
In $d$ dimensions each kernel $\frac{r^m}{r^2}$ becomes $c_d\frac{r^m}{(r^2)^{d/2}}$, with $c_d=\Gamma(\frac d2)\pi^{1-\frac d2}$, the Fourier transform of $-2\pi i\frac{q^m}{q^2}$.
\begin{itemize}
\item $\int d^dx\,K_d(z,z,x)=0$ (scaleless).
\item $\int d^dx\,K_d(y,z,x)=4\pi^2\int\frac{d^dq}{(2\pi)^d}\frac{e^{iq\cdot(y-z)}}{q^2}=\pi^{2-\frac d2}\Gamma(\frac d2-1)[(y-z)^2]^{1-\frac d2}$.
\end{itemize}
With $\big(\frac{\mu^2e^{\gamma_E}}{4\pi}\big)^{\epsilon/2}$,
\begin{equation}
\int d^dx\,\big[K_d(z,z,x)-K_d(y,z,x)\big]=-\pi\,\Gamma\Big(-\frac\epsilon2\Big)\Big(\frac{(y-z)^2\mu^2e^{\gamma_E}}{4}\Big)^{\frac\epsilon2}=\pi\Big[\frac2\epsilon+2\gamma_E+\log\frac{(y-z)^2\mu^2_{\overline{MS}}}{4}\Big]+O(\epsilon).
\label{eq:dimreg}
\end{equation}
The same result follows from the sharp cut $\frac{\theta(|y-z|-|x-z|)}{(x-z)^2}$ and from the smooth form $\frac{1}{(x-z)^2}-\frac{1}{(x-z)^2+(y-z)^2}$ (\texttt{sympy}: @@DR@@).

\subsection*{Step 5: the \texorpdfstring{$p^+$}{p+} integral (exact)}
With $\lambda=\frac{\Lambda}{k^+}$ and $\int_\Lambda^{k^+-\Lambda}\frac{dp^+}{k^+}=\int_\lambda^{1-\lambda}d\xi$:
\begin{equation}
\int_\lambda^{1-\lambda}d\xi\Big[\frac{d_\perp}{d}\xi(1-\xi)+\frac{\xi}{1-\xi}+\frac{1-\xi}{\xi}\Big]=\frac{d_\perp}{d}\Big[\frac16-\lambda^2+\frac23\lambda^3\Big]+2\log\frac{1-\lambda}{\lambda}-2+4\lambda
\end{equation}
(numerically, rel.\ @@PI@@), and $2\log\frac{1-\lambda}{\lambda}=2\log\frac{\vee}{\Lambda}-2\log\frac{\vee}{k^+}+2\log(1-\frac{\Lambda}{k^+})$.

The pole multiplies $\frac{d_\perp}{d}$. With $d=2-\epsilon$,
\begin{equation*}
\frac{d_\perp}{d}\Big[\frac2\epsilon+2\gamma_E+\log\frac{(y-z)^2\mu^2}{4}\Big]=\frac{d_\perp}{2}\Big[\frac2\epsilon+2\gamma_E+\log\frac{(y-z)^2\mu^2}{4}+1\Big]+O(\epsilon)
\end{equation*}
(\texttt{sympy}: @@EXP@@). This is the second line of Eq.~\eqref{eq:UV}.

\subsection*{Step 6: prefactor}
$\frac{g^4}{8\pi^4}\cdot\frac{1}{k^+}\cdot\frac{1}{2\pi}\cdot\pi\cdot N_c=\frac{g^4N_c}{16\pi^4}\frac{1}{k^+}$ (\texttt{sympy}: @@PREF@@). Here:
\begin{itemize}
\item $\frac{1}{2\pi}$ is from $\frac{dp^+}{2\pi}$;
\item one $\frac{1}{k^+}$ of Step 2 turns $dp^+$ into $d\xi$;
\item $\pi$ is from Eq.~\eqref{eq:dimreg};
\item $N_c$ is from the colour factor.
\end{itemize}
Relative to LO, $\frac{1}{(2\pi)^3}\frac{g^2}{2\pi^2k^+}$, the factor is $\frac{\alpha_sN_c}{2\pi}$. The integral is the LO one with the charge at $y$ and the gluon at $z$.

\subsection*{Step 7: the finite term is finite}
At $x\to z$ the $\frac{1}{(x-z)^2}$ terms cancel between 2b and CT for every $p^+$. Numerically, with a smooth test colour factor and the phase kept, $r^2\langle 2b-{\rm CT}\rangle$ at $r=10^{-2}\to10^{-3}$ falls from @@R1@@ to @@R2@@ ($d_\perp=2,3$; $\xi=0.02,0.37,0.95$). Over the same range $r^2\langle 2b\rangle$ stays at $0.09$--$1.75$. At large $|x-z|$ the counterterm falls like $|x-z|^{-3}$, so the finite term behaves there like 2b.

%==========================================================================================
\section{Checks}
Scripts \texttt{part2b\_dperp\_checks.py} (new) and \texttt{part2b\_checks.py} (\texttt{Part2b\_MSbar.pdf}).
\begin{center}\small
\begin{tabular}{@{}p{0.64\textwidth}p{0.3\textwidth}@{}}\toprule
check & result\\\midrule
short-distance coefficient \eqref{eq:coef}, $d=2,3$, $d_\perp=1,2,3,5$ & rel.\ @@SD@@\\
angular average of the counterterm kernel $=\frac{\theta(|y-z|-|x-z|)}{(x-z)^2}$ & @@SAV@@\\
$x$ integral in $d=2-\epsilon$: $K$ form, sharp cut, smooth form (\texttt{sympy}) & @@DR@@\\
$p^+$ integral with $\frac{d_\perp}{d}$ & rel.\ @@PI@@\\
$\frac{d_\perp}{d}$ expansion (\texttt{sympy}) & @@EXP@@\\
$d_\perp=2-\epsilon$ gives Eq.~\eqref{eq:UVcdr} (\texttt{sympy}) & @@CDR@@\\
prefactor (\texttt{sympy}) & @@PREF@@\\
finite term: $r^2\langle2b-{\rm CT}\rangle\to0$ & @@R1@@ $\to$ @@R2@@\\\bottomrule
\end{tabular}
\end{center}
\end{document}
